\section*{Chapter \ref{ch:iv:mock-interviews} --- Mock Interviews}

\begin{solution}{iq:iv:mock-interviews:1}
$0.15 \times 64\,000 = 9\,600$: ten per cent is 6\,400 and five per cent half of that, 3\,200. Any decomposition said aloud scores
(\cref{ch:iv:mental-arithmetic}).
\iqlookfor{speed and a method said aloud.}
\end{solution}

\begin{solution}{iq:iv:mock-interviews:2}
The fifteen remaining tosses have mean 7.5 and standard deviation $\sqrt{15/4} \approx 1.94$, so the total has fair value 10.5 and
the same standard deviation. A market of 10 at 11 is centred and one point wide, about half a standard deviation. The interviewer's
buying carries no information about future tosses, and the three heads are known to both sides, so the market should not move; the
candidate who raises it after a lift has confused flow with information (\cref{ch:iv:betting-and-market-making-games}). Being short
one at 11 against 10.5 is worth 0.5 in expectation.
\iqlookfor{the conditional mean and spread, and no update on uninformative flow.}
\end{solution}

\begin{solution}{iq:iv:mock-interviews:3}
At even money with probability $p = 10/16$ of winning, the Kelly fraction is $2p - 1 = 1/4$, a stake of 100. The expected log growth
per game is $0.625 \ln 1.25 + 0.375 \ln 0.75 \approx 0.0316$. Betting everything maximises the expected bankroll ($1.25 \times 400$
against $1.0625 \times 400$ for the Kelly stake) but loses everything with probability 0.375 each game, so over many games the
bankroll goes to zero almost surely. If the probability is itself uncertain, bet a fraction of Kelly (\cref{ch:iv:betting-and-market-making-games}).
\iqlookfor{the Kelly fraction, the stake, and the difference between the mean and the growth rate.}
\end{solution}

\begin{solution}{iq:iv:mock-interviews:4}
Read as a correlation between signal and next-day return measured on about 2\,500 independent observations, its $t$-statistic is
about $r\sqrt n = 0.02 \times 50 = 1$: not significant on its own. Autocorrelated observations (a slow signal, overlapping returns)
make the effective sample smaller and the $t$ smaller still. A cross-sectional information coefficient of 0.02, computed each day
over many stocks, is a different object: its daily values can be averaged over 2\,500 days and, by the fundamental law, a small
coefficient applied across a broad universe can carry a useful information ratio if turnover and costs allow
(\cref{ch:iv:statistics}).
\iqlookfor{the $t$-statistic, the effective sample size, and breadth.}
\end{solution}

\begin{solution}{iq:iv:mock-interviews:5}
Adding sector dummies changes the question from ``do cheap stocks outperform'' to ``do cheap stocks outperform their own sector''.
The sign flip is an omitted-variable effect: value is correlated with sector membership, and the cheap sectors outperformed over the
sample. Believe the within-sector coefficient for a sector-neutral strategy, the pooled one for a strategy that takes sector bets,
and test the stability of each over subperiods before believing either.
\iqlookfor{omitted-variable bias, and matching the regression to the strategy.}
\end{solution}

\begin{solution}{iq:iv:mock-interviews:6}
$500 \times 10 \times 252 = 1\,260\,000$ stock-days, but only 2\,520 days, and returns on one day share a market shock. Run a
cross-sectional regression each day and test the mean of the 2\,520 daily slopes with a Newey--West standard error, or pool with
standard errors clustered by day. Define the overnight return from the close to the opening auction and the intraday return from the
opening auction to the close, so that both are tradable; exclude stale opens. Add costs (the strategy trades the whole universe
daily), write the rejection criterion before looking, and correct for every variant tried (\cref{ch:iv:research-case-studies-and-take-homes}).
\iqlookfor{the effective sample, clustering, tradable prices, costs and multiple testing.}
\end{solution}

\begin{solution}{iq:iv:mock-interviews:7}
Insert identifiers in a hash set and stop at the first already present: expected $O(n)$ time and $O(n)$ memory. For ten million
identifiers a general-purpose set takes hundreds of megabytes; a sort and a scan of adjacent pairs takes $O(n \log n)$ time and can
run on disk; a Bloom filter gives a fast first pass with false positives only, each checked exactly
(\cref{ch:iv:algorithms-and-data-structures}).
\iqlookfor{the hash set, its memory cost, and alternatives under constraints.}
\end{solution}

\begin{solution}{iq:iv:mock-interviews:8}
A hash map from key to a node of a doubly linked list kept in recency order: \texttt{get} finds the node and moves it to the front,
\texttt{put} inserts or updates at the front and, above capacity, removes the tail node and its map entry; both are $O(1)$. In
Python, \texttt{OrderedDict} with \texttt{move\_to\_end} and \texttt{popitem(last=False)} does the same. The chapter's test runs the
capacity-two sequence (B is evicted) and 200 random sequences against a list-based model.
\iqlookfor{the two structures together, and a test with an eviction.}
\end{solution}

\begin{solution}{iq:iv:mock-interviews:9}
A token bucket of capacity 20 refilled at 100 tokens a second. At time zero 20 messages go and 30 do not. By half a second the bucket
is full again; messages then arrive every 10 milliseconds, exactly the refill rate, so each finds a token and all 100 go: 120 of 150
without waiting, as the chapter's simulation confirms. If the 30 queue instead, they leave one every 10 milliseconds, the last at
0.3 seconds; the bucket is full again by 0.5 seconds, so the second wave is unaffected and all 150 go, 30 of them late. With
a longer first burst or a sooner second wave the queue would delay the later messages too, so the policy (queue, reject, or
give cancels their own budget) must be decided with the desk, who know how quickly a stale order becomes a liability
(\cref{ch:iv:system-design}).
\iqlookfor{the token-bucket model, the count, and the queueing policy.}
\end{solution}

\begin{solution}{iq:iv:mock-interviews:10}
Split by time, never at random: train on the earliest period, validate on the next, test on the last, with a gap (an embargo) at least
as long as the label horizon between them so that overlapping labels do not leak; or walk forward with expanding windows. Hyperparameters
are chosen on validation only, and the test set is used once (\cref{ch:iv:machine-learning}).
\iqlookfor{chronological splits, the embargo, and a test set used once.}
\end{solution}

\begin{solution}{iq:iv:mock-interviews:11}
With one per cent positives, the classifier that always says ``normal'' is 99\% accurate and useless. Report precision and recall
on the anomaly class (or the precision-recall curve and its area), and choose the threshold from the costs of a missed anomaly and a
false alarm.
\iqlookfor{the base rate, and metrics and thresholds tied to costs.}
\end{solution}

\begin{solution}{iq:iv:mock-interviews:12}
$\mathrm{PSI} = \sum_i (a_i - e_i)\ln(a_i/e_i)$ with $e_i = 0.25$ and $a = (0.1, 0.2, 0.3, 0.4)$ gives $0.137 + 0.011 + 0.009 +
0.071 \approx 0.228$, above the common 0.2 threshold for a large shift. First rule out a pipeline fault (units, a changed source, a
stale join); then check whether the model's live performance has moved; retrain only if it has, and on data that include the new
regime (One Quant Book~12, chapter~27).
\iqlookfor{the computation, and diagnosis before retraining.}
\end{solution}

\begin{solution}{iq:iv:mock-interviews:13}
$F = S e^{(r - q)T} = 100 e^{0.04} \approx 104.08$ (\cref{ch:iv:options-and-derivatives}).
\iqlookfor{the carry formula without hesitation.}
\end{solution}

\begin{solution}{iq:iv:mock-interviews:14}
With zero rates, $d_1 = \sigma\sqrt T/2 = 0.1$ and $\varphi(0.1) \approx 0.397$, so vega is $S\varphi(d_1)\sqrt T \approx 39.7$ per
unit of volatility, 0.397 per volatility point. Check: the at-the-money call is about $0.4\,\sigma S\sqrt T$, here 7.97, which is
linear in $\sigma$, so vega is about $7.97/20 \approx 0.40$ per point.
\iqlookfor{the formula, the units, and an independent check.}
\end{solution}

\begin{solution}{iq:iv:mock-interviews:15}
$(C(99) - C(101))/2 \approx 0.4602$, against the digital's $\Phi(d_2) = \Phi(-0.1) \approx 0.4602$ in the model: the centred spread
approximates the derivative of the call price in the strike with an error proportional to the square of the half-width, which the
chapter's test checks at half-widths 4, 2 and 1. In practice the spread's payoff is a ramp, not a step; the desk hedges with a spread
whose payoff dominates the digital's and charges the difference, and the price of the spread reflects the skew, which adds a term
proportional to the slope of implied volatility in the strike (One Quant Book~5, chapter~15).
\iqlookfor{the replication, its error order, the over-hedge and the skew.}
\end{solution}

\begin{solution}{iq:iv:mock-interviews:16}
On the candidate's 24 monthly returns: annualised Sharpe ratio 1.05 (monthly mean over monthly standard deviation, times
$\sqrt{12}$) and maximum drawdown $-2.4\%$ from peak. With two years, the standard error is about $\sqrt{(1 + \mathrm{SR}^2/2)/2}
\approx 0.9$, so the interval runs from about zero to two (\cref{ch:iv:statistics}). The firm recomputes from independent statements.
\iqlookfor{the numbers, and their uncertainty volunteered.}
\end{solution}

\begin{solution}{iq:iv:mock-interviews:17}
Gross return $1.2 \times 8\% = 9.6\%$; costs $20 \times 5$ basis points $= 1\%$; net Sharpe $8.6/8 \approx 1.08$. With impact
growing like the square root of size, ten times the capital multiplies the cost per turn by $\sqrt{10} \approx 3.16$, to about
3.2\%, and the net Sharpe ratio falls to about 0.80. The practical answer is a staged increase with the costs measured at each stage
(One Quant Book~7, chapter~28).
\iqlookfor{costs in Sharpe units, a scaling law for impact, and staging.}
\end{solution}

\begin{solution}{iq:iv:mock-interviews:18}
For $X_t = \mu t + \sigma W_t$, the probability of reaching $-a$ by $T$ is $\Phi\bigl((-a - \mu T)/(\sigma\sqrt T)\bigr) +
e^{-2\mu a/\sigma^2}\,\Phi\bigl((-a + \mu T)/(\sigma\sqrt T)\bigr)$, here $\Phi(-1.75) + e^{-1.5}\Phi(0.25) \approx 0.040 + 0.134
\approx 0.17$, confirmed by simulation (\cref{ch:iv:stochastic-calculus}). That is a loss of 7.5\% from the start. A stop measured
from the running peak is hit far more often: the chapter's simulation of the same process gives about 0.53. The candidate who knows
which rule the firm applies, and what it implies for a manager with a true Sharpe ratio of one, negotiates from facts (One Quant
Book~16, chapter~8).
\iqlookfor{the first-passage formula, the number, and the start-versus-peak distinction.}
\end{solution}
